\documentclass[11pt]{article}
\usepackage[margin=1in]{geometry}
\usepackage{hyperref}
\usepackage{enumitem}

\title{Website Text --- All Pages}
\author{Ibrahim Elsharkawy}
\date{}

\begin{document}
\maketitle
\tableofcontents
\newpage

%% ============================================================
%% HOME PAGE (index.html)
%% ============================================================
\section{Home Page (index.html)}

% ---- Hero / Intro ----
\subsection{Hero Section}
\label{sec:home-hero}

\paragraph{Pretitle}
Ibrahim Elsharkawy

\paragraph{Title}
Physicist \&\\
ML Researcher

\paragraph{Subtitle}
PhD candidate at the University of Toronto. Building foundation models for science at Berkeley Lab/NERSC.

% ---- About Preview ----
\subsection{About Preview}
\label{sec:home-about}

\paragraph{Section Pretitle}
About/CV

\paragraph{Section Title}
Physics for AI,\\AI for Physics.

\paragraph{Bio Paragraph 1}
 My research has been/is focused on using physics-inspired theory and physics data to understand neural network behavior, building foundation models that transfer across scientific domains, and developing AI tools to automate scientific discovery.

\paragraph{Bio Paragraph 2}
I am a PhD candidate at the University of Toronto advised by Prof.Yoni Kahn and doctoral researcher at Berkeley Lab (NERSC) advised by Wahid Bhimji and Benjamin Nachman



% ---- Research Interests Preview ----
\subsection{Research Interests Preview}
\label{sec:home-research}

\paragraph{Section Pretitle}
Research

\paragraph{Section Title}
Research Interests

\paragraph{Lead Text}
We hope to use physics to build better AI, and use AI to do physics better.

\subsubsection{Card 1: Physics Data for Understanding ML}
\label{sec:home-card1}
\textbf{Tag:} Physics for AI\\
\textbf{Title:} Physics Data for Understanding ML\\
\textbf{Description:} We hope to leverage the unique properties of physics datasets (access to controllable data generation and known data symmetry and structure) to probe the internal mechanisms of deep neural networks.

\subsubsection{Card 2: Physics-Inspired Theory for Scaling Laws}
\label{sec:home-card2}
\textbf{Tag:} Physics for AI\\
\textbf{Title:} Physics-Inspired Theory for Scaling Laws\\
\textbf{Description:} Using an effective field theory to predict how neural network ensembles behave, deriving scaling laws for uncertainty quantification without the need to train an ensemble.

\subsubsection{Card 3: Automating Scientific Model Building}
\label{sec:home-card3}
\textbf{Tag:} AI for Physics\\
\textbf{Title:} Automating Scientific Model Building\\
\textbf{Description:} Developing ML methods with the goal of "theory inversion". Parameter estimation with Uncertainty Quantification, Simulation Emulation, and Automating Theory Writing.

\subsubsection{Card 4: Cross-Domain Foundation Models}
\label{sec:home-card4}
\textbf{Tag:} AI for Physics\\
\textbf{Title:} Cross-Domain Foundation Models\\
\textbf{Description:} Building foundation models for scientific point clouds (irregular graphs!) that transfer knowledge across the domains of particle physics, cosmology, and molecular dynamics.

% ---- Photography Preview ----
\subsection{Photography Preview}
\label{sec:home-photography}

\paragraph{Title}
Photography

\paragraph{Lead Text}
Some fun Astrophotography and Portrait work.

% ---- Blog Preview ----
\subsection{Blog Preview}
\label{sec:home-blog}

\paragraph{Section Pretitle}
Blog

\paragraph{Section Title}
Ideas

\paragraph{Coming Soon Text}
 Check back soon.


%% ============================================================
%% RESEARCH PAGE (research.html)
%% ============================================================
\section{Research Page (research.html)}

% ---- Page Header ----
\subsection{Page Header}
\label{sec:research-header}

\paragraph{Pretitle}
Research

\paragraph{Scrolling Marquee Text (repeating loop)}
Physics for AI for Physics for AI for ...

% ---- Research Interest Cards ----
\subsection{Research Interest Cards}
\label{sec:research-cards}

\subsubsection{Card 1: Physics Data for Understanding ML}
\label{sec:research-card1}
\textbf{Tag:} Physics for AI\\
\textbf{Title:} Physics Data for Understanding ML\\
\textbf{Description:} We hope to leverage the unique properties of physics datasets (access to controllable data generation and known data symmetry and structure) to probe the internal mechanisms of deep neural networks.

\subsubsection{Card 2: Physics-Inspired Scaling Laws}
\label{sec:research-card2}
\textbf{Tag:} Physics for AI\\
\textbf{Title:} Physics-Inspired Scaling Laws\\
\textbf{Description:} Using an effective field theory to predict how neural network ensembles behave, deriving scaling laws for uncertainty quantification without the need to train an ensemble.

\subsubsection{Card 3: Automating Scientific Model Building}
\label{sec:research-card3}
\textbf{Tag:} AI for Physics\\
\textbf{Title:} Automating Scientific Model Building\\
\textbf{Description:} Developing ML methods with the goal of ``theory inversion''. Parameter estimation with Uncertainty Quantification, Simulation Emulation, and Automating Theory Writing.

\subsubsection{Card 4: Cross-Domain Foundation Models}
\label{sec:research-card4}
\textbf{Tag:} AI for Physics\\
\textbf{Title:} Cross-Domain Foundation Models\\
\textbf{Description:} Building foundation models for scientific point clouds (irregular graphs!) that transfer knowledge across the domains of particle physics, cosmology, and molecular dynamics.

% ---- Select Papers ----
\subsection{Select Papers}
\label{sec:research-papers}

\begin{enumerate}
    \item \textbf{Uncertainty Quantification from Ensemble Variance Scaling Laws in Deep Neural Networks}\\
    Ibrahim Elsharkawy, et al.\\
    Machine Learning: Science and Technology, vol.\ 6, no.\ 3, 2025

    \item \textbf{Contrastive Normalizing Flows for Uncertainty-Aware Parameter Estimation}\\
    Ibrahim Elsharkawy and Yonatan Kahn\\
    arXiv preprint, 2025

    \item \textbf{OmniMol: Transferring Particle Physics Knowledge to Molecular Dynamics with Point-Edge Transformers}\\
    Ibrahim Elsharkawy, et al.\\
    arXiv preprint, 2026

    \item \textbf{FAIR Universe HiggsML Uncertainty Dataset and Competition}\\
    Wahid Bhimji, et al.\\
    NeurIPS 2025, Dataset and Competition Track

    \item \textbf{OmniCosmos: Transferring Particle Physics Knowledge across the Cosmos}\\
    Vinicius Mikuni, et al.\\
    arXiv preprint, 2025
\end{enumerate}

% ---- Select Talks & Posters ----
\subsection{Select Talks \& Posters}
\label{sec:research-talks}

\begin{itemize}
    \item \textbf{[Invited Talk]} 1st Place Competition Milestone (Ensembles and Uncertainty Quantification)\\
    The Challenge of Handling Uncertainties in Fundamental Science @ NeurIPS 2024

    \item \textbf{[Invited Talk]} Contrastive Normalizing Flows for Uncertainty-Aware Parameter Estimation\\
    CERN 7th Inter-Experimental LHC Machine Learning Workshop, 2025

    \item \textbf{[Poster]} FAIR Universe HiggsML Uncertainty Dataset and Competition\\
    NeurIPS 2025, Dataset and Competition Track

    \item \textbf{[Poster]} Uncertainty Quantification from Scaling Laws in Deep Neural Networks\\
    Machine Learning and the Physical Sciences Workshop @ NeurIPS 2024

    \item \textbf{[Poster]} Pre-Training For Science: A study on Foundation Model Training Objectives\\
    Stanford Center of Decoding the Universe Forum, 2025
\end{itemize}

% ---- CTA ----
\subsection{Call to Action}
\label{sec:research-cta}
Interested in collaborating?


%% ============================================================
%% ABOUT/CV PAGE (about.html)
%% ============================================================
\section{About/CV Page (about.html)}

% ---- Page Header ----
\subsection{Page Header}
\label{sec:about-header}

\paragraph{Pretitle}
About/CV

\paragraph{Title}
Hi, I'm Ibrahim

% ---- Quick Introduction ----
\subsection{A Quick Introduction}
\label{sec:about-intro}

I'm a physicist and machine learning researcher currently pursuing my PhD at the University of Toronto. My work centers on the intersection of physics and artificial intelligence---building foundation models for scientific data at Berkeley Lab/NERSC, developing physics-inspired theories of neural network behavior, and creating AI tools for automated scientific discovery.

% ---- Education ----
\subsection{Education}
\label{sec:about-education}

\begin{itemize}
    \item \textbf{2025--Present:} PhD, Physics --- University of Toronto
    \item \textbf{2023--2025:} MSc, Physics --- University of Illinois Urbana-Champaign
    \item \textbf{2019--2023:} B.S., Computational Physics --- Rice University
    \item \textbf{2019--2023:} B.A., Applied Mathematics \& Philosophy --- Rice University
\end{itemize}

% ---- Experience ----
\subsection{Experience}
\label{sec:about-experience}

\begin{itemize}
    \item \textbf{2025--Present: Doctoral Researcher} --- Berkeley Lab\\
    Building foundation models for scientific point clouds/graphs for High Energy Physics, Cosmology and Molecular Dynamics at NERSC.

    \item \textbf{2024--2025: Research Geophysicist Intern} --- ExxonMobil\\
    ML parameter estimation methods for 3D seismic inversion. Novel ML anomaly detection tools for 4D Seismic.

    \item \textbf{2018--2023: Research Geophysics \& Software Intern} --- Petroleum GeoServices\\
    ML methods for first break picking, seismic anomaly detection, cloud migration.
\end{itemize}

% ---- Awards ----
\subsection{Awards}
\label{sec:about-awards}

\begin{itemize}
    \item \textbf{2025:} University of Toronto Connaught International Fellow
    \item \textbf{2024:} 1st Place Higgs Uncertainty Challenge @ NeurIPS, CERN
    \item \textbf{2023:} 3-Year UIUC Graduate College Fellowship
    \item \textbf{2021:} J \& M Graham Physics Scholar
    \item \textbf{2019:} Shell Eco-Marathon
    \item \textbf{2016:} Eagle Scout
\end{itemize}

% ---- Teaching ----
\subsection{Teaching}
\label{sec:about-teaching}

\begin{itemize}
    \item \textbf{2026:} TA and Lecturer, University of Toronto --- Physics 2108/2109 ``The Physics of Machine Learning''
    \item \textbf{2025:} TA, University of Toronto --- Physics 252 ``Thermal Physics''
    \item \textbf{2024:} TA, UIUC --- Physics 413 ``Optics Lab''
\end{itemize}

% ---- Select Papers (same as research page) ----
\subsection{Select Papers}
\label{sec:about-papers}
(Same list as Research page --- see Section~\ref{sec:research-papers})

% ---- Select Talks (same as research page) ----
\subsection{Select Talks \& Posters}
\label{sec:about-talks}
(Same list as Research page --- see Section~\ref{sec:research-talks})


%% ============================================================
%% CONTACT PAGE (contact.html)
%% ============================================================
\section{Contact Page (contact.html)}

\subsection{Page Header}
\label{sec:contact-header}

\paragraph{Pretitle}
Contact

\paragraph{Title}
Get In Touch

\subsection{Lead Text}
\label{sec:contact-lead}
 I'd like to hear from you!


%% ============================================================
%% BLOG PAGE (blog.html)
%% ============================================================
\section{Blog Page (blog.html)}

\subsection{Page Header}
\label{sec:blog-header}

\paragraph{Pretitle}
Blog

\paragraph{Title}
Ideas

\subsection{Coming Soon}
\label{sec:blog-coming-soon}
Check back soon.


\end{document}
